\documentclass[11pt]{article}
\usepackage[margin=1in]{geometry}
\usepackage{amsmath,amssymb,amsthm}
\usepackage[hidelinks]{hyperref}
\usepackage{xurl}
\newtheorem{lemma}{Lemma}
\newtheorem{theorem}[lemma]{Theorem}
\newtheorem{corollary}[lemma]{Corollary}
\newcommand{\C}{\mathbb C}
\title{Symmetric maximal-braid augmentations and free reflection directions}
\author{Internal source comparison and written proof}
\date{17 November 2026}
\begin{document}
\maketitle

Cornwell's augmentation-rank raising theorem produces examples beyond
the explicit three-braid family in his introduction. Its output has
an additional symmetry condition on the cord values. That condition
allows a direct comparison with the universal reflection obstruction
for the eight knots
$13a_{373},13a_{375},13a_{391},13a_{392},
13a_{3264},13a_{3266},13a_{3275},13a_{3276}$.
The comparison uses a free vector space with a possibly singular
Gram form; it does not require an augmentation to retain its rank
when the meridian parameter is specialized to minus one.

\section{The exact imported algebra and construction}

Use the conventions of \cite[Sections~2.2 and~5.2]{cornwell}.
The free algebra $\mathcal A_n$ has generators $a_{ij}$,
$i\ne j$. Its quotient $\mathcal A_n^{\mathrm{ab}}$ here is
commutative and also imposes $a_{ij}=a_{ji}$; it is not just
the ordinary commutative quotient. For a braid $B\in B_n$,
the extended braid action defines $\Phi_B^L$ by
\[
 \phi_B(a_{i,n+1})=\sum_j(\Phi_B^L)_{ij}a_{j,n+1}.
\]
The Chain Rule is
\[
 \Phi_{BB'}^L=\phi_B(\Phi_{B'}^L)\Phi_B^L.
\]
Write $\Delta(B)=\operatorname{diag}((-1)^{w(B)},1,\ldots,1)$.
Cornwell's Theorem~1.10 and its proof show that a rank-$n$
augmentation of the closure of an $n$-braid satisfies
$\epsilon(\Phi_B^L)=\Delta(B)$.

Theorem~5.4, printed p.~20 of the retained author version,
assumes that the closure of $B\in B_n$ has a rank-$n$
augmentation descending to $\mathcal A_n^{\mathrm{ab}}$.
For $|u|\ge2$, $|v|\ge3$, $\delta=\pm1$ and $1\le i<n$,
it gives a rank-$(n+1)$ augmentation with the same symmetry
for the closure of
\[
 B\sigma_n^\delta\sigma_{i,n}^u\sigma_n^v,
 \qquad
 \sigma_{i,n}=\sigma_i\cdots\sigma_{n-2}\sigma_{n-1}
                 \sigma_{n-2}^{-1}\cdots\sigma_i^{-1},
\]
provided $u+v-1$ is even, or $i=1$ and $u$ is odd.
This is a genuine higher-braid family; its bridge-four instances
must be compared separately from Theorem~1.13's three-braid cases.
The full augmentation-rank raising proof is imported, rather than
reproved here.

\section{A transport that does not divide by a Gram determinant}

\begin{lemma}\label{lem:transport}
Suppose a map $\epsilon:\mathcal A_n^{\mathrm{ab}}\to\C$
satisfies $\epsilon(\Phi_B^L)=\Delta(B)$.
Then the group of the closed braid has a representation on a free
$n$-dimensional vector space, preserving a symmetric form, which
sends the $n$ strand meridians to nonisotropic orthogonal reflections
with linearly independent normal vectors. The form may be singular.
\end{lemma}
\begin{proof}
Let $V$ have basis $e_1,\ldots,e_n$, and give it the symmetric form
\[
 G(e_i,e_i)=2,\qquad G(e_i,e_j)=\epsilon(a_{ij}).
\]
For a vector of norm two write $R_u(x)=x-G(u,x)u$.
Its displacement image is $\C u$, even when $G$ is singular,
because $(R_u-I)u=-2u$.

Assign $R_{e_i}$ to the initial strand meridians. At a positive
braid generator $\sigma_k$, propagate root vectors by
\[
 (u_k,u_{k+1})\longmapsto
       (R_{u_k}u_{k+1},u_k).
\]
Conjugation of reflections makes this exactly the Artin action
on the corresponding meridian images. The inverse move is
$(u_k,u_{k+1})\mapsto(u_{k+1},R_{u_{k+1}}u_k)$.
If $g=G(u_k,u_{k+1})$, the positive root coefficient block is
\[
                 \begin{pmatrix}-g&1\\1&0\end{pmatrix}.
\]
This is the $k,k+1$ block of $\epsilon(\Phi_{\sigma_k}^L)$
in the source conventions. Directly pairing the new vectors gives
\[
 \begin{split}
  g'_{k,k+1}&=-g_{k,k+1},\\
  g'_{k,i}&=g_{k+1,i}-g_{k+1,k}g_{k,i},\\
  g'_{k+1,i}&=g_{k,i}\qquad(i\ne k,k+1),
 \end{split}
\]
with symmetry and all other pairings unchanged. These are precisely
the displayed $\phi_{\sigma_k}$ rules of Section~2.2 after imposing
$a_{ij}=a_{ji}$. The inverse move gives the inverse rules.
Consequently the coefficient matrix for the final root vectors
satisfies the same Chain Rule as $\Phi_B^L$, and is
$\epsilon(\Phi_B^L)$.

By hypothesis each final vector is its initial basis vector up
to the diagonal sign in $\Delta(B)$. A sign does not change a
reflection. All closed-braid relations therefore hold, yielding
the required knot- or link-group representation.
The initial vectors are independent by the definition of $V$;
at no step have we inverted $G$ or quotiented out its radical.
\end{proof}

\section{The eight hyperbolic alternating knots}

Let
\[
 \mathcal S=\{13a_{373},13a_{375},13a_{391},13a_{392},
13a_{3264},13a_{3266},13a_{3275},13a_{3276}\}.
\]
The separate written proof
\path{UNIVERSAL_REFLECTION_OBSTRUCTION.tex} shows that in any
orthogonal reflection representation of a knot in $\mathcal S$,
over characteristic different from two, its four specified
generating meridian normals span a space of dimension at most three.
This statement includes singular forms and arbitrary ambient
dimensions. It is an explicit dependency of the following theorem.

\begin{theorem}\label{thm:exclude}
For $K\in\mathcal S$ and every $n$-braid $B$ whose closure is $K$,
there is no map $\epsilon:\mathcal A_n^{\mathrm{ab}}\to\C$
with $\epsilon(\Phi_B^L)=\Delta(B)$ when $n\ge4$.
In particular no maximal-rank augmentation of such a braid closure
descends to $\mathcal A_n^{\mathrm{ab}}$.
\end{theorem}
\begin{proof}
If a map existed, Lemma~\ref{lem:transport} would give a representation
with independent strand-meridian normals $e_1,\ldots,e_n$.
Let $U$ be the span of the four specified meridian normals from
the universal obstruction. It has dimension at most three.
Every generating reflection preserves $U$ and displaces every
ambient vector by an element of $U$. The same holds for every
word in them. Since those four meridians generate the knot group,
it holds in particular for each strand-meridian reflection $R_{e_i}$.
But $(R_{e_i}-I)e_i=-2e_i$, so every $e_i$ belongs to $U$.
Thus $n\le\dim U\le3$, a contradiction.

For the augmentation assertion, apply the necessary identity from
Cornwell's Theorem~1.10 to its symmetric cord values. This step
uses the augmentation at its original parameter, not a claim about
its rank at minus one.
\end{proof}

\begin{corollary}\label{cor:raising}
None of the eight knots in $\mathcal S$ is a bridge-four output of
Cornwell's Theorem~5.4. More generally it admits no full-rank
four-braid augmentation descending to the source's symmetric
quotient $\mathcal A_4^{\mathrm{ab}}$.
\end{corollary}

The symmetry hypothesis is essential to the comparison. The theorem
does not bound the ordinary augmentation rank of these knots, exclude
nonsymmetric rank-four augmentations, or exclude rank-four
augmentations presented on more than four strands. No worldwide
priority claim follows from this family exclusion.

\section{Exact scope of the replay and source audit}

The retained primary source is Cornwell's author version
arXiv:1310.7526v3, dated 5 June 2017,
\path{sources/cornwell-kch-1310.7526v3.pdf}.
Section~2.2 supplies the braid action and Chain Rule;
Theorem~1.10's proof is on printed p.~18; the symmetric quotient
is defined on p.~19; Theorem~5.4 is on p.~20 and proved on p.~23.
The primary formulas and these exact theorem hypotheses were read.
The original construction theorem and knot-contact-homology
correspondence are imports, not independently established by this
comparison.

The standard-library checker
\path{verify_symmetric_augmentation_transport.py} verifies every
positive and negative local move on a generic symmetric four-root
Gram matrix, both inverse cancellations, the adjacent braid relation,
and disjoint-generator commutation. It also verifies independent
displacement directions for a rank-one Gram matrix, illustrating
why a rank drop of a specialized Gram form is immaterial.
It does not rerun the eight knot-specific universal obstructions;
those have their own checker and retained inputs. The all-braid
transport and the exclusion are the written arguments above,
supported by these finite algebra controls, not formal verification
or external specialist review.

\begin{thebibliography}{9}
\bibitem{cornwell} C. R. Cornwell,
KCH representations, augmentations, and $A$-polynomials,
Journal of Symplectic Geometry 15 (2017), 983--1017.
Author version arXiv:1310.7526v3, Sections~2.2 and~5.2--5.3,
Theorems~1.10 and~5.4.
\url{https://arxiv.org/abs/1310.7526v3}.
\end{thebibliography}
\end{document}
